\section{Limitations}

The thesis uses a controlled experimental setup, but the conclusions are still
bounded by the dataset, model search, and evaluation design. The most important
limitations are:
\begin{itemize}
	\item the experiments use a 2,000-file subset rather than the full available corpus;
	\item training files were sampled at the MIDI-file level rather than selecting one MIDI representation per unique music IDs, which can reduce musical diversity when several MIDI versions represent the same song;
	\item the Transformer models are matched baselines, not independently optimized architectures;
	\item the xLSTM hyperparameter search is broader than the Transformer search, so the comparison is controlled but not equally tuned for all model families;
	\item the medium xLSTM follow-up changes capacity, training length, and schedule together, so its overfitting behavior does not isolate one cause;
	\item the context length is fixed at 1024 tokens, which can be too short for dense or complex musical passages where one pattern already occupies a large part of the available window;
	\item the generation protocol uses one shared stochastic sampling setup rather than a separate sampling search per model;
	\item automatic metrics do not fully measure listening-level properties such as phrasing, repetition, harmonic direction, or long-range coherence;
	\item the hearing-test material is partial and informal, not a blind listening study;
	\item generation-time measurements are practical run observations, not controlled hardware benchmarks;
	\item the Gradio applications support qualitative inspection but do not by themselves provide user-study evidence.
\end{itemize}

The dataset is the first limiting factor. The 2,000-file subset makes the
experiments feasible on rented and local GPUs, but it is small relative to the
full LMD-derived corpus. The subset is also based on MIDI files rather than a
music-ID-first selection, so songs with multiple MIDI versions can occupy more
of the training set than musically distinct songs. Genre tags are LastFM-derived,
multi-label, and noisy, so they are useful for reporting and split diagnostics
but not for clean genre-conditioned conclusions. The split construction reduces
leakage by using music IDs and normalized note signatures, but it cannot
guarantee that every cover version, transposition, alternate arrangement, or
metadata mismatch is removed.

The model comparison is also limited by the amount of tuning applied to each
family. The xLSTM setup is derived from several earlier validation sweeps,
whereas GPT-2 and LLaMA are trained as controlled baselines around the selected
xLSTM setting. This makes the comparison useful for asking how the architectures
behave under the same representation and training budget. It does not prove that
the chosen Transformer sizes, schedules, or regularization settings are the best
possible Transformer configurations for this dataset.

The fixed context length also limits what the models can learn from complex
pieces. A 1024-token window may cover a useful amount of music for simpler
passages, but dense arrangements, fast rhythmic activity, ornamentation, or many
simultaneous events can consume that window quickly. In those cases the model
may see only a short local fragment of a larger motif or form-level pattern,
making it effectively short-sighted even when the original MIDI contains
long-range structure.

The generation evaluation should be read as a structured proxy for musical
quality. MusPy descriptors and FMD provide reproducible measurements, but they
cannot fully capture whether a continuation has convincing phrasing, musical
development, harmonic motion, or arrangement-level coherence. The partial WAV
renders help with manual listening, but they were prepared as inspection
material rather than as a controlled listening test with randomized samples,
raters, rating criteria, and statistical analysis.

There are additional limitations in the generation protocol itself. All models
use the same prompt construction and sampling settings, which keeps the
comparison fair but may not be optimal for any individual model. The
\verb|generation_length: reference| setting also produces continuations with
reference-matched lengths instead of a fixed duration or a hard token cap. This
is useful for paired metric computation, but it can create very long generations
and complicates both runtime comparison and audio rendering.

Finally, the implementation artifacts are useful but not a complete deployment
study. The Gradio applications show that the trained checkpoints can be loaded,
sampled from, continued, and inspected interactively. They do not measure user
preference, latency under concurrent use, robustness to arbitrary uploaded MIDI
files, or the usability of the interface for musicians. Those questions require
a separate evaluation beyond the scope of the current thesis.
